% !TEX root = ../main.tex
\corrigesChapitre{13}{Ensembles et raisonnement}
\corrigesSouples

\begin{blocLivre}{cor-13-1}
\begin{corrige}{1}{Appartenance et opérations}
\demarche{Séparer les éléments communs, les éléments de la réunion et les éléments extérieurs.}
\begin{questions}
\item $\res{A\cap B=\{2\}}$, $\res{A\cup B=\{1;2;3;4;6\}}$ et $\res{E\setminus A=\{1;3;5\}}$.
\item On écrit $\res{2\in A}$ et $\res{\{2\}\subseteq A}$. Le nombre 2 est un élément, le singleton $\{2\}$ est un sous-ensemble.
\end{questions}
\label{figlivre-r13-01-venn-cor}
\figureLivre[corrige]{r13-01-venn}

\end{corrige}
\end{blocLivre}
\begin{blocLivre}{cor-13-2}
\begin{corrige}{2}{Réunir des intervalles}
\demarche{Comparer les bornes et vérifier, pour chacune, si elle est incluse.}
\begin{questions}
\item Les valeurs communes donnent $\res{]1;3]}$. La réunion couvre sans trou $\res{[-2;5[}$.
\item $\res{\R\setminus A=\mathopen{]}-\infty;-2[\cup]3;+\infty[}$. Les bornes $-2$ et 3 sont exclues du complément car elles appartiennent à A.
\end{questions}
\label{figlivre-r13-02-intervalles-cor}
\figureLivre[corrige]{r13-02-intervalles}

\end{corrige}
\end{blocLivre}
\begin{blocLivre}{cor-13-3}
\begin{corrige}{3}{Construire des couples}
\demarche{Respecter l’ordre des deux coordonnées d’un couple.}
\begin{questions}
\item $A\times B=\{(a;0),(a;1),(a;2),(b;0),(b;1),(b;2)\}$, de cardinal $\res{6}$.
\item $(1;a)$ n’appartient pas à $A\times B$, car sa première coordonnée n’est pas dans A. Il appartient à $B\times A$.
\end{questions}
\label{figlivre-r13-03-couples-cor}
\figureLivre[corrige]{r13-03-couples}

\end{corrige}
\end{blocLivre}
\begin{blocLivre}{cor-13-4}
\begin{corrige}{4}{Trouver un contre-exemple}
\demarche{Un seul exemple qui contredit une affirmation universelle suffit à la réfuter.}
\begin{questions}
\item Pour $x=1/2$, $x^2=1/4<1/2$ : l’affirmation est fausse.
\item $x=-2$ vérifie $x^2=4$, mais pas $x=2$. Les deux racines de $x^2=4$ sont $-2$ et 2.
\end{questions}
\end{corrige}
\end{blocLivre}
\begin{blocLivre}{cor-13-5}
\begin{corrige}{5}{Nier une phrase quantifiée}
\demarche{Échanger les quantificateurs et nier la propriété avec la bonne inégalité.}
\begin{questions}
\item La négation est « il existe un réel $x$ tel que $x^2<0$ ». La phrase initiale est vraie, sa négation fausse.
\item La négation est « pour tout réel $x$, $x^2\neq-1$ ». La phrase initiale est fausse, sa négation vraie.
\end{questions}
\end{corrige}
\end{blocLivre}
\begin{blocLivre}{cor-13-6}
\begin{corrige}{6}{Réciproque et contraposée}
\demarche{Distinguer l’implication, sa réciproque et sa contraposée.}
\begin{questions}
\item Pour $x>2$, $x+2>0$ et $x-2>0$, donc $x^2-4=(x-2)(x+2)>0$.
\item La réciproque est « si $x^2>4$, alors $x>2$ ». Elle est fausse pour $x=-3$.
\item La contraposée est « si $x^2\leq4$, alors $x\leq2$ ». Elle est vraie, car équivalente à l’implication initiale.
\end{questions}
\end{corrige}
\end{blocLivre}
\begin{blocLivre}{cor-13-7}
\begin{corrige}{7}{Nécessaire ou suffisante}
\demarche{Identifier quel énoncé suffit à garantir l’autre.}
\begin{questions}
\item $P\Rightarrow Q$ : si $n=4k$, alors $n=2(2k)$ est pair.
\item P est \textbf{suffisante} pour Q ; Q est \textbf{nécessaire} pour P.
\item $n=6$ est pair sans être multiple de 4. Q n’est donc pas suffisante pour P.
\end{questions}
\end{corrige}
\end{blocLivre}
\begin{blocLivre}{cor-13-8}
\begin{corrige}{8}{Variable, inconnue et paramètre}
\demarche{Le rôle de la lettre dépend de la tâche : décrire, résoudre ou choisir une fonction.}
\begin{questions}
\item $x$ est une variable ; l’identité est vraie pour tout réel $x$.
\item $x$ est une inconnue dont on cherche les valeurs : $\res{x=-1\text{ ou }1}$.
\item Pour une fonction fixée de la famille, $x$ est la variable et $m$ un paramètre constant. Faire varier $m$ change la fonction étudiée.
\end{questions}
\end{corrige}
\end{blocLivre}
\begin{blocLivre}{cor-13-9}
\begin{corrige}{9}{Complémentaires et connecteurs}
\demarche{Passer de la négation d’un ou à la conjonction de deux négations.}
\begin{questions}
\item $A\cup B=\{1;2;3;4;6;8\}$, donc le complémentaire est $\res{\{5;7\}}$.
\item $E\setminus A=\{1;3;5;7\}$ et $E\setminus B=\{5;6;7;8\}$ ; l’intersection est aussi $\res{\{5;7\}}$.
\item Ne pas être « dans A ou dans B » signifie être « hors de A et hors de B ». C’est une loi de De Morgan.
\end{questions}
\label{figlivre-r13-09-venn-cor}
\figureLivre[corrige]{r13-09-venn}

\end{corrige}
\end{blocLivre}
\begin{blocLivre}{cor-13-10}
\begin{corrige}{10}{Lire une quantification implicite}
\demarche{Pour réfuter une croissance sur un domaine, exhiber deux points dans le mauvais ordre.}
\begin{questions}
\item Il existe $a,b\in I$ avec $a<b$ et $f(a)>f(b)$.
\item Prendre $a=-2$ et $b=-1$ : $a<b$ mais $f(a)=4>1=f(b)$.
\item Pour $0\leq a<b$, $b^2-a^2=(b-a)(b+a)>0$. Donc $f(a)<f(b)$ : la fonction y est même strictement croissante.
\end{questions}
\end{corrige}
\end{blocLivre}
\begin{blocLivre}{cor-13-11}
\begin{corrige}{11}{Démontrer par disjonction des cas}
\demarche{Traiter tous les cas possibles, puis formuler une conclusion commune.}
\begin{questions}
\item Si $n$ est pair, $n=2k$ et le produit est multiple de 2. Si $n$ est impair, $n+1$ est pair, donc le produit est encore multiple de 2. Les deux cas couvrent tous les entiers naturels.
\item Si $x\geq0$, $|x|=x$. Si $x<0$, $|x|=-x>x$. Ainsi l’inégalité est vraie dans les deux cas.
\end{questions}
\end{corrige}
\end{blocLivre}
\begin{blocLivre}{cor-13-12}
\begin{corrige}{12}{Une preuve par contraposée}
\demarche{Prouver la contraposée revient à prouver l’implication de départ.}
\begin{questions}
\item La contraposée est : « si un entier est pair, alors son carré est pair ».
\item Si $n=2k$, $n^2=4k^2=2(2k^2)$ est pair. La contraposée étant vraie, l’implication initiale l’est aussi.
\item Si $n=2k+1$, $n^2=4k^2+4k+1=2(2k^2+2k)+1$ est impair. Ainsi $\res{n^2\text{ impair}\iff n\text{ impair}}$.
\end{questions}
\end{corrige}
\end{blocLivre}
\begin{blocLivre}{cor-13-13}
\begin{corrige}{13}{Une preuve par l'absurde}
\demarche{Supposer une fraction irréductible, puis obtenir un facteur commun impossible.}
\begin{questions}
\item En élevant $\sqrt2=p/q$ au carré, on obtient $p^2=2q^2$ : le carré $p^2$ est pair.
Si $p$ était impair, son carré serait impair d’après la propriété précédente.
Donc $p$ est pair et s’écrit $p=2k$ pour un entier positif $k$.
\item En remplaçant $p$, $4k^2=2q^2$, donc $q^2=2k^2$.
Le même argument impose que $q$ soit pair.
\item Alors 2 divise simultanément $p$ et $q$.
Cela contredit le choix initial d’une fraction irréductible.
L’hypothèse $\sqrt2=p/q$ est donc impossible, et $\res{\sqrt2\notin\mathbb Q}$.
\end{questions}

\end{corrige}
\end{blocLivre}
\begin{blocLivre}{cor-13-14}
\begin{corrige}{14}{Une transformation sous condition}
\demarche{Conserver la valeur interdite pendant toutes les transformations de l’équation.}
\begin{questions}
\item Le domaine est $\R\setminus\{1\}$. Sur ce domaine, $(x^2-1)/(x-1)=x+1$.
\item $x+1=2$ donne $x=1$, qui est interdit. L’ensemble des solutions est $\res{\varnothing}$.
\item La multiplication conduit à $(x-1)^2=0$, mais cette équation n’est équivalente à l’originale que pour $x\neq1$. Oublier cette condition ajoute une fausse solution.
\end{questions}
\end{corrige}
\end{blocLivre}
\begin{blocLivre}{cor-13-15}
\begin{corrige}{15}{Quantifier une famille de trinômes}
\demarche{Fixer le paramètre et comparer le minimum du trinôme à zéro.}
\begin{questions}
\item $x^2-2mx+1=(x-m)^2+1-m^2$.
Pour chaque $m$ fixé, le carré est minimal en $x=m$ : le minimum de $P_m$ vaut $1-m^2$.
\item La condition pour tout réel $x$ impose que ce minimum soit positif ou nul :
$1-m^2\geq0\iff m^2\leq1\iff\res{m\in[-1;1]}$.
\item Pour être strictement positif en tout réel, le minimum doit être strictement positif :
$1-m^2>0\iff\res{m\in\mathopen{]}-1;1[}$.
\item La négation est : il existe un réel $x$ tel que $P_m(x)<0$.
Elle se produit pour $\res{m<-1\text{ ou }m>1}$ : le choix explicite $x=m$ donne alors $P_m(m)=1-m^2<0$.
\item Pour $m=1$ ou $m=-1$, le minimum vaut 0 et est réellement atteint en $x=m$.
Ces deux paramètres satisfont la condition large, mais pas la condition stricte.
L’explication orale doit donc distinguer positivité du minimum et possibilité d’une valeur nulle.
\end{questions}

\end{corrige}
\end{blocLivre}
